\chapter{Getallenreeksen}\label{ch:b1:series}

Oneindig veel getallen optellen betekent de limiet van de partiële
sommen nemen --- niets meer en niets minder. Dit hoofdstuk zet de
definities op en de convergentiecriteria die in het eerste jaar
bruikbaar zijn: vergelijking en equivalenten voor positieve termen,
het \omterm{thm:b1:series:ratio}{quotiëntcriterium}, de vergelijking met een \omterm{thm:b1:integration:def}{integraal} die de
\omterm{thm:b1:series:riemann}{Riemann-reeksen} levert, de \omterm{thm:b1:series:absolute}{absolute convergentie} en de stelling over
\omterm{thm:b1:series:alternating}{alternerende reeksen}. De fijnere theorie (producten van \omterm{def:b1:series:def}{reeksen},
sommatie per pakket, \omterm{def:b1:series:def}{reeksen} van functies) hoort bij het tweede
jaar.

\section{Algemeenheden}

\begin{definition}\label{def:b1:series:def}
Gegeven een rij $(u_n)$ is de \emph{reeks}\index{reeks} $\sum u_n$
de rij van \emph{partiële sommen} $S_N = \sum_{n=0}^{N} u_n$. De
reeks \emph{convergeert} wanneer $(S_N)$ convergeert; de limiet is
de \emph{som} $\sum_{n=0}^{\infty} u_n$, en $R_N = \sum_{n > N} u_n
= S - S_N$ is de \emph{restterm}, die naar $0$ nadert.
\end{definition}

\begin{example}[Meetkundige reeks]\label{ex:b1:series:geometric}
Voor $q \in \C$ is $\;S_N = \sum_{n=0}^{N} q^n = \frac{1 -
q^{N+1}}{1-q}$ ($q \neq 1$). De \omterm{def:b1:series:def}{reeks} convergeert dan en slechts dan
als $\abs q < 1$ (\cref{exo:b1:seq:3}), met\index{meetkundige reeks}
\[
\sum_{n=0}^{\infty} q^n = \frac{1}{1 - q} .
\]
\end{example}

\begin{example}[Periodieke decimalen zijn meetkundige reeksen]\label{ex:b1:series:periodicdecimal}
Welk getal is $0.363636\dots$? Zijn schrijfwijze zelf is een
\omterm{def:b1:series:def}{reeks}:
\[
0.\overline{36} = \sum_{k=1}^{\infty} \frac{36}{100^k}
= 36\cdot\frac{1/100}{1 - 1/100} = \frac{36}{99} =
\frac{4}{11} ,
\]
via de meetkundige som met $q = \frac{1}{100}$. In het algemeen
is een blok $B$ van $p$ cijfers dat zich eeuwig herhaalt
$\frac{B}{10^p - 1}$ waard --- het mechanisme achter het
periodiciteitscriterium van \cref{pb:b1:reals:1}, dat de taal van
dit hoofdstuk eindelijk in één regel formuleert: een \omterm{pb:b1:reals:1}{decimale
ontwikkeling} is een convergente \omterm{def:b1:series:def}{reeks}, uiteindelijk periodiek
precies wanneer haar som rationaal is. De cijfermachinerie van
hoofdstuk 10, daar met kale \omterm{def:b1:reals:bounds}{suprema} gebouwd, was
reeksentheorie in vermomming.
\end{example}

\begin{proposition}[Eerste feiten]\label{prop:b1:series:first}
\begin{enumerate}
  \item Convergeert $\sum u_n$, dan is $u_n \to 0$. (Het omgekeerde
        is \emph{onwaar}: de harmonische \omterm{def:b1:series:def}{reeks}.)
  \item Lineariteit: convergente \omterm{def:b1:series:def}{reeksen} mag men optellen en
        schalen, met de verwachte sommen.
  \item (Telescoperen\index{telescoperende reeks}) $\sum (v_{n+1} -
        v_n)$ convergeert dan en slechts dan als $(v_n)$
        convergeert, met som $\lim v_n - v_0$.
  \item Eindig veel termen wijzigen verandert de convergentie niet
        (alleen de som).
\end{enumerate}
\end{proposition}

\begin{proof}
(1) $u_N = S_N - S_{N-1} \to S - S = 0$. De harmonische \omterm{def:b1:series:def}{reeks} heeft
$u_n = \frac1n \to 0$ en divergeert toch (\cref{exo:b1:seq:5}). (2)
Bewerkingen met limieten. (3) $S_N = v_{N+1} - v_0$. (4) De partiële
sommen veranderen met een uiteindelijk constante hoeveelheid.
\end{proof}

\begin{example}[Cijfers plannen met de meetkundige restterm]\label{ex:b1:series:georemainder}
Voor $\abs q < 1$ is de restterm van de \omterm{ex:b1:series:geometric}{meetkundige reeks}
expliciet:
\[
R_N = \sum_{n = N+1}^{\infty} q^n = \frac{q^{N+1}}{1 - q} .
\]
Dit zet nauwkeurigheidsdoelen om in aantallen termen nog vóór
enige berekening. Om $\sum_{n\geq0} \bigl(\frac13\bigr)^n =
\frac32$ tot op $10^{-10}$ te evalueren: we hebben
$\frac{(1/3)^{N+1}}{2/3} \leq 10^{-10}$ nodig, dat wil zeggen
$3^{N} \geq \frac{3}{2}\cdot 10^{10}$, dus $N \geq 22$ (want
$3^{22} \approx 3.1\cdot10^{10}$): drieëntwintig termen, op
voorhand bekend. Elke schatting met meetkundige snelheid uit de
weekendopgaven (de $\frac13$-reeks voor $\ln 2$, de boogtangensen
van Machin in \cref{pb:b1:taylor:1}) is dit budget van twee
regels in beroepskledij.
\end{example}

\begin{example}[Een langere telescoop]\label{ex:b1:series:telescope3}
Bereken $\sum_{n\geq1} \frac{1}{n(n+1)(n+2)}$. Splitsing in
partieelbreuken (\cref{ch:b1:fractions}):
\[
\frac{1}{n(n+1)(n+2)}
= \frac{1/2}{n} - \frac{1}{n+1} + \frac{1/2}{n+2}
= \frac12\Bigl(\frac{1}{n(n+1)} - \frac{1}{(n+1)(n+2)}\Bigr),
\]
waarbij de tweede vorm --- een verschil van opeenvolgende waarden
van $w_n = \frac{1}{n(n+1)}$ --- de telescoperende is. Bijgevolg
\[
\sum_{n=1}^{N} \frac{1}{n(n+1)(n+2)}
= \frac12\Bigl(w_1 - w_{N+1}\Bigr)
= \frac12\Bigl(\frac12 - \frac{1}{(N+1)(N+2)}\Bigr)
\longrightarrow \frac14 .
\]
Het afsluitende inzicht: partieelbreuken met drie termen
telescoperen zelden zoals ze er staan; hergroepeer ze eerst tot
een verschil $w_n - w_{n+1}$ --- de beloning is niet alleen
convergentie maar de exacte som, die geen enkel
vergelijkingscriterium ooit levert.
\end{example}

\section{Reeksen met niet-negatieve termen}

\begin{theorem}[Begrensde partiële sommen]\label{thm:b1:series:positive}
Is $u_n \geq 0$ voor alle $n$, dan stijgen de partiële sommen, dus:
$\sum u_n$ convergeert $\iff$ haar partiële sommen zijn naar boven
begrensd. Bijgevolg het \emph{vergelijkingscriterium}: geldt $0 \leq
u_n \leq v_n$ voor alle (grote) $n$, dan
\[
\sum v_n \text{ convergeert} \implies \sum u_n \text{ convergeert},
\qquad
\sum u_n \text{ divergeert} \implies \sum v_n \text{ divergeert}.
\]
En het \emph{equivalentiecriterium}: is $u_n \sim v_n$ met $v_n \geq
0$, dan hebben de twee \omterm{def:b1:series:def}{reeksen} dezelfde aard.
\end{theorem}

\begin{proof}
De stelling over monotone limieten (\cref{thm:b1:seq:monotone}) voor
het eerste punt; vergelijking van partiële sommen voor het tweede.
Equivalenten: voor grote $n$ is $\frac12 v_n \leq u_n \leq 2 v_n$
(definitie van $\sim$ met $\varepsilon = \frac12$), en de
vergelijking werkt beide kanten op.
\end{proof}

\begin{example}[Een equivalent dat divergentie bewijst]\label{ex:b1:series:equivdiverge}
Wat is de aard van $\sum_{n\geq1} n\sin\dfrac{1}{n^2}$? Omdat
$\frac{1}{n^2} \to 0$ en $\sin h \sim h$ in $0$:
\[
n\sin\frac{1}{n^2} \;\sim\; n\cdot\frac{1}{n^2} = \frac1n ,
\]
en het equivalentiecriterium draagt de divergentie van de
harmonische \omterm{def:b1:series:def}{reeks} over: divergent --- ook al naderen de termen tot
$0$. Eén ontwikkeling, één schaal, één oordeel; hetzelfde patroon
in twee stappen (equivalent, daarna opzoeken bij Riemann of bij
de \omterm{ex:b1:series:geometric}{meetkundige reeks}) beslist alle vier de \omterm{def:b1:series:def}{reeksen} van
\cref{exo:b1:series:3}.
\end{example}

\begin{example}[Het equivalentiecriterium in één regel]\label{ex:b1:series:equivtest}
Wat is de aard van $\sum_{n \geq 1} \frac{\sqrt{n+1} - \sqrt
n}{n}$? Maak de teller toegevoegd:
\[
\frac{\sqrt{n+1} - \sqrt n}{n}
= \frac{1}{n\,(\sqrt{n+1} + \sqrt n)}
\sim \frac{1}{2\,n^{3/2}} ,
\]
een convergente Riemann-schaal ($\alpha = \frac32 > 1$): de \omterm{def:b1:series:def}{reeks}
convergeert. De hele beslissing kostte één equivalent en één
opzoeking --- mits de termen niet-negatief zijn, wat hier het
geval is. Het afsluitende inzicht: voor positieve \omterm{def:b1:series:def}{reeksen} is de
volledige convergentietheorie een \emph{woordenboek van schalen}
($n^{-\alpha}$, $q^n$, $\frac{1}{n(\ln n)^\alpha}$) plus de
vrijheid om een term door een equivalent te vervangen; het
analytische werk zit in de asymptotiek (\cref{ch:b1:taylor}),
nooit in het sommeren.
\end{example}

\begin{theorem}[Vergelijking met een integraal; Riemann-reeksen]\label{thm:b1:series:riemann}
Zij $f$ \omterm{def:b1:continuity:continuous}{continu}, niet-negatief en \emph{dalend} op
$\intco{1}{+\infty}$. Dan geldt
\[
\int_1^{N+1} f(t)\,\dd t \;\leq\; \sum_{n=1}^{N} f(n) \;\leq\; f(1)
+ \int_1^{N} f(t)\,\dd t ,
\]
dus convergeert $\sum f(n)$ dan en slechts dan als $\bigl(\int_1^x
f\bigr)$ begrensd is. In het bijzonder geldt voor $\alpha \in
\R$:\index{Riemann-reeks}
\[
\sum_{n \geq 1} \frac{1}{n^\alpha} \text{ convergeert}
\iff \alpha > 1,
\]
en $\sum_{n=1}^{N} \frac1n = \ln N + O(1)$.
\end{theorem}

\begin{proof}
Voor $n \leq t \leq n+1$ geeft de monotonie $f(n+1) \leq f(t)
\leq f(n)$; integratie over $\intcc{n}{n+1}$ (een segment van
lengte $1$) levert
\[
f(n+1) \;\leq\; \int_n^{n+1} f(t)\,\dd t \;\leq\; f(n) .
\]
Sommatie van de rechterongelijkheden voor $n = 1, \dots, N-1$
geeft $\int_1^{N} f \leq \sum_{n=1}^{N-1} f(n)$, en dus de
bovenste omsluiting na optelling van $f(N) \leq f(1)$; sommatie
van de linkerongelijkheden voor $n = 1, \dots, N$ geeft
$\sum_{n=2}^{N+1} f(n) \leq \int_1^{N+1} f$, wat na herindexering
de onderste omsluiting is. Convergentie: de partiële sommen en de
\omterm{thm:b1:integration:def}{integralen} $\int_1^x f$ begrenzen elkaar op de constante $f(1)$
na, en beide zijn stijgend, zodat de ene begrensd is dan en
slechts dan als de andere dat is
(\cref{thm:b1:series:positive}). Voor $f(t) = t^{-\alpha}$
($\alpha \neq 1$) is $\int_1^x t^{-\alpha}\dd t =
\frac{x^{1-\alpha} - 1}{1 - \alpha}$, begrensd dan en slechts dan
als $\alpha > 1$; voor $\alpha = 1$ is de \omterm{thm:b1:integration:def}{integraal} $\ln x \to
\infty$, en geeft de omsluiting $\ln(N+1) \leq H_N \leq 1 + \ln
N$. Voor $\alpha \leq 0$ naderen de termen niet tot $0$.
\end{proof}

\begin{example}[De harmonische stapel]\label{ex:b1:series:harmonicstack}
Hoeveel termen moet de harmonische \omterm{def:b1:series:def}{reeks} opstapelen om $20$ te
overschrijden? De omsluiting $\ln(N+1) \leq H_N \leq 1 + \ln N$
antwoordt zonder ook maar iets te sommeren: $H_N \geq 20$ vereist
$1 + \ln N \geq 20$, dus $N \geq \eu^{19} \approx 1.8\cdot10^{8}$,
en is gegarandeerd zodra $\ln(N + 1) \geq 20$, dus $N \approx
\eu^{20} \approx 4.9\cdot10^{8}$. (De weekendopgave scherpt dit
aan tot $N \approx \eu^{20 - \gamma} \approx 2.7\cdot10^{8}$ via
de \omterm{pb:b1:series:1}{constante van Euler}.) Het afsluitende inzicht: de vergelijking
met een \omterm{thm:b1:integration:def}{integraal} beslist niet alleen over convergentie --- zij
\emph{lokaliseert} partiële sommen met logaritmische
nauwkeurigheid, en maakt van een hopeloze berekening (honderden
miljoenen termen) een schatting van twee regels.
\end{example}

\begin{theorem}[Quotiëntcriterium (d'Alembert)]\label{thm:b1:series:ratio}
Zij $u_n > 0$ met $\frac{u_{n+1}}{u_n} \to
\ell$.\index{quotiëntcriterium}
\begin{itemize}
  \item Is $\ell < 1$, dan convergeert $\sum u_n$;
  \item is $\ell > 1$, dan is $u_n \to +\infty$: divergentie;
  \item is $\ell = 1$, dan is er geen besluit ($\sum \frac1n$
        divergeert, $\sum \frac{1}{n^2}$ convergeert).
\end{itemize}
\end{theorem}

\begin{proof}
Is $\ell < 1$, leg dan $q \in \intoo{\ell}{1}$ vast: voorbij zekere
$N$ is $u_{n+1} \leq q\,u_n$, dus $u_n \leq u_N q^{\,n-N}$ met
inductie: vergelijking met een \omterm{ex:b1:series:geometric}{meetkundige reeks}. Is $\ell > 1$, dan
is de rij $(u_n)$ voorbij zekere $N$ stijgend, zodat zij niet naar
$0$ kan naderen (haar limiet is, als die bestaat, $\geq u_N > 0$);
volgens \cref{prop:b1:series:first}~(1) volgt divergentie --- en in
feite geeft $u_n \geq u_N q^{n-N}$ met $q > 1$ dat $u_n \to \infty$.
\end{proof}

\begin{example}\label{ex:b1:series:ratio}
$\sum \frac{x^n}{n!}$ convergeert voor elke $x > 0$: het quotiënt
is $\frac{x}{n+1} \to 0$. Haar som is $\eu^x$: volgens
Taylor--Lagrange (\cref{thm:b1:taylor:lagrange}) op $\intcc{0}{x}$
is
\[
\Bigl| \eu^x - \sum_{k=0}^{n} \frac{x^k}{k!} \Bigr|
\leq \eu^{x}\, \frac{x^{n+1}}{(n+1)!} \xrightarrow[n\to\infty]{} 0 ,
\]
waarbij de grens naar $0$ nadert omdat de faculteit domineert
(\cref{exo:b1:integration:9}~(1) gebruikte hetzelfde feit).
Hetzelfde argument sommeert de \omterm{def:b1:series:def}{reeksen} van $\sin$, $\cos$, $\sinh$
en $\cosh$ op heel $\R$.
\end{example}

\begin{example}[Het quotiëntcriterium is voldoende, niet
noodzakelijk]\label{ex:b1:series:rationotnecessary}
Zij $u_n = 2^{-n}$ voor even $n$ en $u_n = 2^{-n-2}$ voor oneven
$n$. De opeenvolgende quotiënten schommelen tussen $\frac{1}{8}$
en $\frac12\cdot4 = 2$, zodat $\frac{u_{n+1}}{u_n}$ geen limiet
heeft en d'Alembert zwijgt --- terwijl $u_n \leq 2^{-n}$ en het
vergelijkingscriterium de convergentie ogenblikkelijk beslecht.
De hypothese van het criterium (het quotiënt \emph{convergeert})
is een echte beperking: zij past bij termen met één dominante
multiplicatieve structuur (faculteiten, machten), en faalt op
alles wat ademt. Gedragen de quotiënten zich slecht, stap dan
terug naar vergelijking met een meetkundige omhulling --- en dat
is alles wat het \omterm{thm:b1:series:ratio}{quotiëntcriterium} ooit was, zoals zijn bewijs
laat zien.
\end{example}

\begin{example}[Het quotiëntcriterium in faculteitengevechten]\label{ex:b1:series:ratiobinom}
Wat is de aard van $\sum_{n\geq0} \dfrac{(n!)^2}{(2n)!}$ (de
omgekeerden van de centrale \omterm{def:b1:counting:objects}{binomiaalcoëfficiënten}, op de factor
$n + 1$ na)? Het quotiënt laat de faculteiten instorten:
\[
\frac{u_{n+1}}{u_n}
= \frac{((n+1)!)^2}{(n!)^2}\cdot\frac{(2n)!}{(2n+2)!}
= \frac{(n+1)^2}{(2n+1)(2n+2)}
\longrightarrow \frac14 < 1 :
\]
convergent, met ruimte over --- de termen nemen in wezen af als
$4^{-n}$, in overeenstemming met $\binom{2n}{n} \geq
\frac{4^n}{2n+1}$ uit \cref{pb:b1:integration:1}. Het
afsluitende inzicht: quotiënten van faculteiten zijn precies wat
het \omterm{thm:b1:series:ratio}{quotiëntcriterium} verteert --- elke faculteit valt weg tot
een rationale functie van $n$, waarvan de limiet uit de leidende
termen wordt afgelezen.
\end{example}

\section{Absolute convergentie; alternerende reeksen}

\begin{theorem}[Absolute convergentie]\label{thm:b1:series:absolute}
Convergeert $\sum \abs{u_n}$ (\emph{absolute
convergentie}\index{absolute convergentie}), dan convergeert $\sum
u_n$, en is $\bigl|\sum u_n\bigr| \leq \sum \abs{u_n}$. Dit geldt
voor reële zowel als complexe termen.
\end{theorem}

\begin{proof}
De partiële sommen voldoen, voor $M > N$ (criterium van Cauchy,
\cref{thm:b1:seq:complete}), aan
\[
\abs{S_M - S_N} = \Bigl| \sum_{n=N+1}^{M} u_n \Bigr|
\leq \sum_{n=N+1}^{M} \abs{u_n},
\]
wat klein is voor grote $N$, omdat de partiële sommen van
$\sum\abs{u_n}$ een \omterm{def:b1:seq:cauchy}{Cauchyrij} vormen. Dus is $(S_N)$ een \omterm{def:b1:seq:cauchy}{Cauchyrij}
en bijgevolg convergent. De ongelijkheid gaat over op de limiet
vanuit de eindige driehoeksongelijkheid.
\end{proof}

\begin{example}[Absolute convergentie, reëel en complex]\label{ex:b1:series:absexamples}
$\sum_{n\geq1} \frac{\sin n}{n^2}$: de termen wisselen grillig
van teken (de rij $(\sin n)$ ligt zelfs \omterm{def:b1:topology:dense}{dicht} in
$\intcc{-1}{1}$, \cref{exo:b1:seq:12}), en er is geen
alternerende structuur in zicht. De \omterm{thm:b1:series:absolute}{absolute convergentie} redt
alles in één klap: $\bigl|\frac{\sin n}{n^2}\bigr| \leq
\frac{1}{n^2}$, een convergente schaal, dus convergeert de \omterm{def:b1:series:def}{reeks}.
Hetzelfde schild werkt over $\C$: $\sum_{n\geq1}\frac{\eu^{\iu
n}}{n^2}$ convergeert omdat $\bigl|\frac{\eu^{\iu n}}{n^2}\bigr|
= \frac{1}{n^2}$ --- tekenpatronen, zelfs tweedimensionale, doen
er niet toe zodra de moduli sommeerbaar zijn. Het afsluitende
inzicht: de \omterm{thm:b1:series:absolute}{absolute convergentie} is het enige gereedschap van
dit hoofdstuk dat nooit vraagt hoe de tekens zijn georganiseerd;
probeer haar eerst (\cref{met:b1:series:decide}), en bewaar de
delicate criteria voor de \omterm{def:b1:series:def}{reeksen} die erop stuklopen.
\end{example}

\begin{theorem}[Criterium voor alternerende reeksen]\label{thm:b1:series:alternating}
Zij $(a_n)$ dalend met $a_n \to 0$. Dan convergeert de \omterm{thm:b1:series:alternating}{alternerende
reeks} $\sum (-1)^n a_n$; haar som ligt tussen elke twee
opeenvolgende partiële sommen, en\index{alternerende reeks}
\[
\abs{R_N} = \Bigl| \sum_{n > N} (-1)^n a_n \Bigr| \leq a_{N+1} .
\]
\end{theorem}

\begin{proof}
De even en de oneven partiële sommen zijn ingesloten: $S_{2p+2} -
S_{2p} = a_{2p+2} - a_{2p+1} \leq 0$ (dalend), $S_{2p+1} - S_{2p-1}
= a_{2p} - a_{2p+1} \geq 0$ (stijgend), en $S_{2p} - S_{2p+1} =
a_{2p+1} \to 0$. Volgens \cref{thm:b1:seq:adjacent} delen zij een
limiet $S$, die het criterium van de twee \omterm{def:b1:seq:subsequence}{deelrijen}
(\cref{prop:b1:seq:subsequences}) tot de limiet van $(S_N)$ maakt;
bovendien zit $S$ gevangen tussen opeenvolgende partiële sommen, en
is $\abs{S - S_N}$ hoogstens de opening naar de volgende, $a_{N+1}$.
\end{proof}

\begin{example}[Alternerende harmonische reeks]\label{ex:b1:series:ln2}
$\sum_{n \geq 1} \frac{(-1)^{n-1}}{n}$ convergeert (criterium voor
\omterm{thm:b1:series:alternating}{alternerende reeksen}), maar niet absoluut (harmonische \omterm{def:b1:series:def}{reeks}).
Haar som is $\ln 2$: integreer de eindige meetkundige identiteit
$\frac{1}{1+t} = \sum_{k=0}^{n-1} (-t)^k + \frac{(-t)^n}{1+t}$
over $\intcc{0}{1}$:
\[
\ln 2 = \sum_{k=1}^{n} \frac{(-1)^{k-1}}{k}
+ (-1)^n \int_0^1 \frac{t^n}{1+t}\,\dd t ,
\qquad
0 \leq \int_0^1 \frac{t^n}{1+t}\,\dd t \leq \frac{1}{n+1} \to 0 .
\]
De convergentie is pijnlijk traag ($R_N \approx \frac{1}{N}$) ---
\omterm{thm:b1:series:alternating}{alternerende reeksen} convergeren door opheffing, niet door
kleinheid.
\end{example}

\begin{omfigure}
\begin{tikzpicture}[x=0.78cm, y=7.2cm]
  \draw[-Stealth] (0.3,0.42) -- (11,0.42)
    node[below, font=\small] {$N$};
  \draw (1,0.415) -- (1,0.425) node[below=1pt, font=\small]
    {$1$};
  \draw (5,0.415) -- (5,0.425) node[below=1pt, font=\small]
    {$5$};
  \draw (10,0.415) -- (10,0.425) node[below=1pt, font=\small]
    {$10$};
  \draw[gray, dashed] (0.3,0.6931) -- (10.7,0.6931)
    node[right, font=\small, text=black] {$\ln 2$};
  \node[left, font=\small] at (0.3,1) {$1$};
  \node[left, font=\small] at (0.3,0.5) {$0.5$};
  \draw[omDef, thick]
    plot coordinates {(1,1) (2,0.5) (3,0.8333) (4,0.5833)
    (5,0.7833) (6,0.6167) (7,0.7595) (8,0.6345)
    (9,0.7456) (10,0.6456)};
  \fill[omDef] (1,1) circle (1.5pt);
  \fill[omDef] (2,0.5) circle (1.5pt);
  \fill[omDef] (3,0.8333) circle (1.5pt);
  \fill[omDef] (4,0.5833) circle (1.5pt);
  \fill[omDef] (5,0.7833) circle (1.5pt);
  \fill[omDef] (6,0.6167) circle (1.5pt);
  \fill[omDef] (7,0.7595) circle (1.5pt);
  \fill[omDef] (8,0.6345) circle (1.5pt);
  \fill[omDef] (9,0.7456) circle (1.5pt);
  \fill[omDef] (10,0.6456) circle (1.5pt);
\end{tikzpicture}

{\small De partiële sommen $S_N$ van de alternerende harmonische
\omterm{def:b1:series:def}{reeks} $1 - \frac12 + \frac13 - \cdots$ springen bij elke stap over
hun limiet $\ln 2$ heen: de oneven sommen van boven, de even van
onder, elke sprong ter grootte $\frac{1}{N+1}$. De insluiting is
het bewijs van \cref{thm:b1:series:alternating} zichtbaar gemaakt
--- en het trage sluiten van de tang ($\abs{S_N - \ln 2} \approx
\frac{1}{2N}$, weekendopgave \cref{pb:b1:series:1}) is de reden
waarom niemand $\ln 2$ zo berekent.}
\end{omfigure}

\begin{remark}[Veelgemaakte fouten met reeksen]
(i) \emph{Het equivalentiecriterium heeft een teken nodig}: zij
$v_n = \frac{(-1)^n}{\sqrt n}$ en $u_n = v_n + \frac1n$. Dan is
$\frac{u_n}{v_n} = 1 + \frac{(-1)^n}{\sqrt n} \to 1$, dus $u_n
\sim v_n$; en toch convergeert $\sum v_n$ (criterium voor
\omterm{thm:b1:series:alternating}{alternerende reeksen}) terwijl $\sum u_n = \sum v_n + \sum
\frac1n$ divergeert. Equivalentie beheerst de \emph{grootte} van
de termen, en voor \omterm{def:b1:series:def}{reeksen} met wisselende tekens is grootte geen
lot --- het criterium is geformuleerd, en waar, uitsluitend voor
(uiteindelijk) niet-negatieve termen. (ii) \emph{$u_n \to 0$
bewijst niets}: de harmonische \omterm{def:b1:series:def}{reeks} is het eeuwige
tegenvoorbeeld; de omgekeerde richting
(\cref{prop:b1:series:first}~(1)) is alleen een snel
divergentiecriterium. (iii) \emph{Quotiëntlimiet $1$ is
stilzwijgen, geen convergentie}: zowel $\sum\frac1n$ als
$\sum\frac{1}{n^2}$ heeft quotiënt $\to 1$; stap over op
Riemann-schalen of op vergelijking met een \omterm{thm:b1:integration:def}{integraal}. (iv)
\emph{Alternerend vereist dalend}: $\sum \frac{(-1)^n}{n +
(-1)^n}$ lijkt alternerend en wordt alleen met een ontwikkeling
behandeld (\cref{exo:b1:series:5}); de weekendopgave van
\cref{ch:b1:taylor} (vraag 23 aldaar) toont dat het criterium
zonder monotonie ronduit kan falen. (v) \emph{Groeperen en
herordenen zijn niet gratis}: haakjes invoegen is onschuldig
voor convergente \omterm{def:b1:series:def}{reeksen} maar kan uit divergentie convergentie
scheppen ($1 - 1 + 1 - \cdots$ paarsgewijs gegroepeerd), en
herordenen kan de som zelf veranderen --- het drama dat de
weekendopgave van dit hoofdstuk opvoert
(\cref{pb:b1:series:1}).
\end{remark}

\begin{method}[De aard van een reeks bepalen]\label{met:b1:series:decide}
\begin{enumerate}
  \item Nadert $u_n$ tot $0$? Zo niet: divergentie, stop.
  \item Niet-negatieve termen: zoek een equivalent van $u_n$
        (ontwikkelingen, \cref{ch:b1:taylor}!) en vergelijk met
        Riemann-schalen of meetkundige schalen; faculteiten en
        machten roepen om het \omterm{thm:b1:series:ratio}{quotiëntcriterium}; een dalende
        $f(n)$ roept om vergelijking met een \omterm{thm:b1:integration:def}{integraal}.
  \item Wisselende tekens: probeer eerst de \omterm{thm:b1:series:absolute}{absolute convergentie};
        faalt die, dan het criterium voor \omterm{thm:b1:series:alternating}{alternerende reeksen}
        (controleer \emph{dalend} zorgvuldig); daarbuiten:
        gereedschap van het tweede jaar.
\end{enumerate}
\end{method}

\begin{example}[Oneven noemers, de halve telescoop]\label{ex:b1:series:oddtelescope}
Bereken $\sum_{n \geq 1} \dfrac{1}{4n^2 - 1}$. Partieelbreuken:
$\frac{1}{(2n-1)(2n+1)} = \frac12\bigl(\frac{1}{2n-1} -
\frac{1}{2n+1}\bigr)$, dus
\[
\sum_{n=1}^{N} \frac{1}{4n^2 - 1}
= \frac12\Bigl(1 - \frac{1}{2N+1}\Bigr)
\longrightarrow \frac12 .
\]
Vergelijk met $\sum \frac{1}{n(n+1)} = 1$
(\cref{exo:b1:series:1}): hetzelfde telescoperende skelet, maar
de opeenvolgende termen liggen hier twee uit elkaar in de oneven
getallen, en de factor $\frac12$ registreert de stap. Het
afsluitende inzicht: telescoperen is een verandering van
gezichtspunt, geen truc --- zodra de algemene term een verschil
$w_n - w_{n+1}$ is van een rij met een limiet, is de som $w_1 -
\lim w$, precies \cref{prop:b1:series:first}~(3).
\end{example}

\begin{remark}[De pijplijn van de analyse, terugblikkend]
Dit hoofdstuk is de plek waar de analyse van het volume
samenkomt, en elk criterium noemt zijn voorouder. Begrensde
partiële sommen is de stelling over monotone limieten
(\cref{ch:b1:seq}), zelf het volledigheidsaxioma van
\cref{ch:b1:reals}; de \omterm{thm:b1:series:absolute}{absolute convergentie} is het criterium
van Cauchy; het integraalcriterium is de omsluiting van
oppervlakten uit \cref{ch:b1:integration}; equivalenten van
algemene termen zijn de ontwikkelingen van \cref{ch:b1:taylor};
en de stelling over \omterm{thm:b1:series:alternating}{alternerende reeksen} is het lemma van de
\omterm{thm:b1:seq:adjacent}{ingesloten rijen} in zondagse kleren. Achterstevoren gelezen legt
de pijplijn uit waar elk hoofdstuk \emph{voor} diende --- en de
weekendopgaven die erdoorheen zijn geregen ($b$-adische cijfers,
Cesàro--Stolz, de irrationaliteitsmachines, de \omterm{pb:b1:series:1}{constante van
Euler}) zijn dezelfde paar ideeën die elkaar op steeds grotere
hoogte ontmoeten. De lineaire algebra die volgt verandert van
onderwerp, niet van maatstaven: de gewoonte van exacte
\omterm{def:b1:logic:statement}{uitspraken} met gecertificeerde fout overleeft de overgang van
limieten naar dimensies.
\end{remark}

\begin{remark}[Waar reeksen hierna heen gaan]
Dit hoofdstuk sluit de analyse van het volume af en opent drie
deuren. In het volume van bachelorjaar 2 krijgen \omterm{def:b1:series:def}{reeksen} een
veranderlijke ($\sum a_n x^n$: machtreeksen, met hun
convergentiestraal) en daarna een functiewaardige theorie
(fourierreeksen); de tweedeling absoluut tegenover voorwaardelijk
convergent, gedramatiseerd in de weekendopgave hieronder, wordt de
hoeksteen van beide. In de kansrekening (het volume van
bachelorjaar 3) \emph{zijn} verwachtingswaarden van discrete
stochastische variabelen \omterm{def:b1:series:def}{reeksen}, en is de \omterm{thm:b1:series:absolute}{absolute convergentie}
wat hen welgedefinieerd maakt. En de \omterm{thm:b1:series:riemann}{Riemann-reeks} $\sum n^{-s}$,
doorgetrokken naar complexe $s$, wordt de zètafunctie --- de
meest bestudeerde \omterm{def:b1:series:def}{reeks} van de wiskunde.
\end{remark}

\section{Oefeningen}

\begin{exercise}[$\star$]\label{exo:b1:series:1}
Aard (en som, waar telescoperend) van:
\[
\sum_{n\geq1} \frac{1}{n(n+1)},
\qquad
\sum_{n\geq2} \ln\Bigl(1 - \frac{1}{n^2}\Bigr),
\qquad
\sum_{n\geq0} \frac{3^n + 4^n}{5^n} .
\]
\end{exercise}

\begin{exercise}[$\star$]\label{exo:b1:series:2}
Aard van: $\;\sum \dfrac{n^2}{2^n}$; $\;\sum \dfrac{n!}{n^n}$;
$\;\sum \dfrac{2^n\,n!}{n^n}$; $\;\sum \dfrac{3^n\,n!}{n^n}$.
\emph{(\omterm{thm:b1:series:ratio}{Quotiëntcriterium}; denk aan $\bigl(1 + \frac1n\bigr)^n \to
\eu$.)}
\end{exercise}

\begin{exercise}[$\star$]\label{exo:b1:series:3}
Aard van: $\;\sum \sin\dfrac{1}{n^2}$; $\;\sum \Bigl(1 -
\cos\dfrac1n\Bigr)$; $\;\sum \dfrac{1}{\sqrt{n(n+1)}}$; $\;\sum
\dfrac{\ln n}{n^2}$ \emph{(vergelijk met $n^{-3/2}$)}.
\end{exercise}

\begin{exercise}[$\star$]\label{exo:b1:series:4}
Bewijs dat $\sum_{n\geq1} \frac{1}{n^2}$ convergeert met som $\leq
2$, met behulp van $\frac{1}{n^2} \leq \frac{1}{n(n-1)}$ voor $n
\geq 2$ en een telescoperende grens.
\end{exercise}

\begin{exercise}[$\star\star$]\label{exo:b1:series:5}
Aard van $\;\sum \dfrac{(-1)^n}{\sqrt n}$, van $\;\sum
\dfrac{(-1)^n}{n + (-1)^n}$ \emph{(ontwikkel: het criterium voor
\omterm{thm:b1:series:alternating}{alternerende reeksen} is niet rechtstreeks van toepassing ---
waarom?)}, en van $\;\sum \sin\bigl(\pi\sqrt{n^2+1}\,\bigr)$
\emph{(herleid modulo $\pi$: $\sqrt{n^2+1} = n + \frac{1}{2n} +
O(n^{-3})$)}.
\end{exercise}

\begin{exercise}[$\star\star$]\label{exo:b1:series:6}
Voor welke $\alpha > 0$ convergeert $\sum_{n \geq 2} \dfrac{1}{n
(\ln n)^{\alpha}}$? \emph{(Vergelijking met een \omterm{thm:b1:integration:def}{integraal};
substitueer $u = \ln t$.)}
\end{exercise}

\begin{exercise}[$\star\star$]\label{exo:b1:series:7}
Zij $u_n = \dfrac{1}{n} - \ln\Bigl(1 + \dfrac1n\Bigr)$. Bewijs dat
$0 \leq u_n \leq \dfrac{1}{2n^2}$, dat $\sum u_n$ convergeert, en
leid het bestaan af van de \omterm{pb:b1:series:1}{constante van
Euler}\index{constante van Euler}:
\[
\gamma = \lim_{N \to \infty} \Bigl( \sum_{n=1}^{N} \frac 1n - \ln N
\Bigr) .
\]
\end{exercise}

\begin{exercise}[$\star\star$]\label{exo:b1:series:8}
Bereken de sommen
\[
\sum_{n=1}^{\infty} \frac{1}{n(n+2)}
\qquad\text{en}\qquad
\sum_{n=0}^{\infty} \frac{n}{2^n} .
\]
\emph{(Voor de eerste: partieelbreuken. Voor de tweede: bereken
$\sum_{n=1}^{N} n x^{n-1}$ in gesloten vorm en laat $N \to \infty$
bij $x = \frac12$.)}
\end{exercise}

\begin{exercise}[$\star\star\star$]\label{exo:b1:series:9}
(Condensatie van Cauchy) Zij $(u_n)$ niet-negatief en dalend.
Bewijs dat
\[
\sum_{n \geq 1} u_n \text{ convergeert}
\iff
\sum_{k \geq 0} 2^k\, u_{2^k} \text{ convergeert},
\]
door pakketten termen tussen opeenvolgende machten van $2$ te
vergelijken. Vind daaruit het criterium van Riemann en
\cref{exo:b1:series:6} terug.
\end{exercise}

\begin{exercise}[$\star\star\star$]\label{exo:b1:series:10}
Bewijs met de integraalidentiteit van \cref{ex:b1:series:ln2},
aangepast aan $\frac{1}{1+t^2}$, de formule van Leibniz
\[
\frac{\pi}{4} = \sum_{n=0}^{\infty} \frac{(-1)^n}{2n+1}
= 1 - \frac13 + \frac15 - \frac17 + \cdots
\]
met de foutgrens $\abs{R_N} \leq \frac{1}{2N+3}$.
\end{exercise}

\begin{exercise}[$\star\star$]\label{exo:b1:series:11}
Aard van $\displaystyle\sum_{n \geq 1} \frac{1}{n^{1 + 1/n}}$.
\emph{(Bereken de limiet van $n^{1/n}$ en bepaal een equivalent van
de algemene term: het criterium van Riemann heeft een
\emph{vaste} exponent nodig.)}
\end{exercise}

\begin{exercise}[$\star\star\star$]\label{exo:b1:series:12}
Zij $(u_n)$ niet-negatief en \emph{dalend} met $\sum u_n$
convergent. Bewijs dat $n\,u_n \to 0$ \emph{(begrens $n\,u_{2n}$
door een schijf $\sum_{k=n+1}^{2n} u_k$ en gebruik het criterium
van Cauchy)}. Toon aan dat het omgekeerde faalt en dat de
hypothese van monotonie niet kan worden weggelaten.
\end{exercise}

\section{Opgave: de constante van Euler en de reeks die haar som
verandert}

\begin{problem}[{Weekendopgave --- $H_n = \ln n + \gamma +
\frac{1}{2n} + O(n^{-2})$, en het herschikken van $1 - \frac12 +
\frac13 - \dots$ tot $\frac{\ln 2}{2}$}]
\label{pb:b1:series:1}
Twee verhalen delen de harmonische \omterm{def:b1:series:def}{reeks}. Ten eerste de exacte
boekhouding van haar divergentie: $H_n - \ln n$ convergeert naar de
\omterm{pb:b1:series:1}{constante van Euler} $\gamma$ (\cref{exo:b1:series:7}), en deze
opgave scherpt die \omterm{def:b1:logic:statement}{uitspraak} aan tot een tweezijdige wet
$\frac{1}{2(n+1)} \leq H_n - \ln n - \gamma \leq \frac{1}{2n}$,
die $\gamma = 0.5772\dots$ met de hand certificeert. Ten tweede
het schandaal van de voorwaardelijke convergentie: de alternerende
harmonische \omterm{def:b1:series:def}{reeks} sommeert tot $\ln 2$ (\cref{ex:b1:series:ln2}),
en toch sommeren \emph{dezelfde termen, in een andere volgorde},
tot $\frac{\ln 2}{2}$ --- of tot $\ln 2 + \frac12\ln\frac pq$
voor willekeurige $p, q$, of tot elk reëel getal hoegenaamd
(Riemann). De twee verhalen zijn er één: de herschikte sommen
worden \emph{met} de $\gamma$-wet
berekend.\index{constante van Euler}\index{herschikking van
reeksen}

\medskip
\noindent\textbf{Deel I --- $\gamma$, ingesloten.} Stel $a_n = H_n
- \ln n$ en $b_n = H_n - \ln(n+1)$.
\begin{enumerate}
  \item Toon met $\frac{t}{1+t} \leq \ln(1 + t) \leq t$ aan dat
        $(a_n)$ daalt, $(b_n)$ stijgt, en dat zij ingesloten zijn;
        hun gemeenschappelijke limiet is $\gamma$, met $b_n \leq
        \gamma \leq a_n$ voor elke $n$.
  \item Een eerste numeriek schot: sluit $\gamma$, uitgaande van
        $H_{10} = 2.928968\dots$, in tussen $b_{10} = 0.5311$ en
        $a_{10} = 0.6264$. Hoe groot zou $n$ moeten zijn opdat deze
        grove insluiting vier decimalen geeft?
  \item Toon de exacte staartvoorstelling $a_n - \gamma = \sum_{k
        \geq n} w_k$ aan (als limiet van partiële sommen), waarbij
        \[
        w_k = a_k - a_{k+1}
        = \ln\Bigl(1 + \frac1k\Bigr) - \frac{1}{k+1}
        = \int_k^{k+1} \frac{(k + 1 - t)}{t\,(k+1)}\,\dd t ,
        \]
        en leid uit de integraalvorm de tweezijdige grens
        $\dfrac{1}{2(k+1)^2} \leq w_k \leq \dfrac{1}{2k(k+1)}$ af.
\end{enumerate}

\medskip
\noindent\textbf{Deel II --- De $\frac{1}{2n}$-wet.}
\begin{enumerate}[resume]
  \item Sommeer de grenzen van vraag 3 (beide leden telescoperen of
        laten zich met telescopen vergelijken) en besluit tot de
        wet:
        \[
        \frac{1}{2(n+1)} \;\leq\; H_n - \ln n - \gamma \;\leq\;
        \frac{1}{2n} \qquad (n \geq 1).
        \]
  \item Leid af dat $H_n = \ln n + \gamma + \frac{1}{2n} +
        O\bigl(\frac{1}{n^2}\bigr)$; toon preciezer aan dat
        $\gamma_n = H_n - \ln n - \frac{1}{2n}$ voldoet aan
        $-\frac{1}{2n(n+1)} \leq \gamma_n - \gamma \leq 0$.
  \item Certificeer vier decimalen met $n = 100$: bereken,
        uitgaande van $H_{100} = 5.1873775\dots$, dat
        $\gamma_{100} = 0.577207\dots$ en besluit dat $\gamma =
        0.5772 \pm 5\cdot10^{-5}$ (werkelijke waarde
        $0.5772156\dots$).
  \item Twee dividenden van de wet, die beide later nodig zijn:
        als $m \to \infty$,
        \[
        H_{2m} - H_m = \ln 2 - \frac{1}{4m} +
        O\Bigl(\frac{1}{m^2}\Bigr),
        \qquad
        \sum_{j=1}^{m} \frac{1}{2j-1} = \frac{\ln m}{2} + \ln 2
        + \frac\gamma2 + o(1) ,
        \]
        de tweede via $\sum_{j \leq m} \frac{1}{2j-1} = H_{2m} -
        \frac12 H_m$, en evenzo $\sum_{j=1}^{m} \frac{1}{2j} =
        \frac{\ln m}{2} + \frac\gamma2 + o(1)$.
\end{enumerate}

\medskip
\noindent\textbf{Deel III --- De alternerende harmonische \omterm{def:b1:series:def}{reeks},
tot op tweede orde.}
\begin{enumerate}[resume]
  \item Toon (met inductie of door groeperen) de identiteit
        $\sum_{k=1}^{2m} \frac{(-1)^{k-1}}{k} = H_{2m} - H_m$ aan,
        en leid daaruit zowel de som $\ln 2$ (opnieuw) als de
        exacte snelheid af:
        \[
        \sum_{k=1}^{2m} \frac{(-1)^{k-1}}{k}
        = \ln 2 - \frac{1}{4m} + O\Bigl(\frac{1}{m^2}\Bigr) .
        \]
  \item Leid de asymptotische fout van de alternerende harmonische
        \omterm{def:b1:series:def}{reeks} af bij \emph{elke} index: $S - S_N \sim
        \frac{(-1)^N}{2N}$ --- twee keer kleiner dan de grens voor
        het slechtste geval $a_{N+1} \approx \frac1N$ uit
        \cref{thm:b1:series:alternating}.
  \item (Versnelling voor niets) Toon aan dat de gemiddelde sommen
        $\tilde S_N = \frac{S_N + S_{N+1}}{2}$ voldoen aan $\tilde
        S_N = \ln 2 + O\bigl(\frac{1}{N^2}\bigr)$. Controle:
        $S_{10} = 0.64563$, $S_{11} = 0.73654$, $\tilde S_{10} =
        0.69109$, tegenover $\ln 2 = 0.69315$: één gemiddelde
        koopt twee decimalen.
  \item Leg in twee zinnen uit waarom zo'n truc niets kan
        uitrichten tegen een \emph{positief} verschijnsel als de
        insluiting van vraag 2: de alternerende fout schommelt
        (teken $(-1)^N$), zodat middelen haar leidende term
        opheft, terwijl de fout $\frac{1}{2n}$ van de
        $\gamma$-insluiting een vast teken heeft. ($a_n$ en $b_n$
        middelen \emph{helpt} wel: breng $\frac{a_n + b_n}{2}$ in
        verband met de middenpuntsschatting $H_n - \ln\bigl(n +
        \frac12\bigr)$ en toon aan dat haar fout
        $O\bigl(\frac{1}{n^2}\bigr)$ is.)
\end{enumerate}

\medskip
\noindent\textbf{Deel IV --- Starheid en haar falen.}
\begin{enumerate}[resume]
  \item Toon aan dat het positieve deel $\sum \frac{1}{2j-1}$ en
        het negatieve deel $\sum \frac{1}{2j}$ van de alternerende
        harmonische \omterm{def:b1:series:def}{reeks} beide divergeren --- het kenmerk van
        \emph{voorwaardelijke} convergentie.
  \item Bewijs de algemene \omterm{def:b1:logic:statement}{uitspraak} achter vraag 12: convergeert
        $\sum u_n$ terwijl $\sum \abs{u_n}$ divergeert, dan
        divergeren de \omterm{def:b1:series:def}{reeks} van de positieve delen $\sum u_n^+$ en
        die van de negatieve delen $\sum u_n^-$ beide \emph{(uit
        $u_n^\pm = \frac{\abs{u_n} \pm u_n}{2}$: convergeerde er
        één, dan ook de andere, en dus $\sum\abs{u_n}$)}. Dit
        onuitputtelijke reservoir van positieve en negatieve massa
        is wat het recept van Riemann zal uitgeven.
  \item (Starheid) Bewijs: convergeert $\sum u_n$ \emph{absoluut}
        en is $\sigma \colon \N \to \N$ een bijectie, dan
        convergeert $\sum u_{\sigma(n)}$ naar dezelfde som
        \emph{(voor grote $N$ bevatten de eerste $M$ herschikte
        termen $u_0, \dots, u_N$; vergelijk de partiële sommen via
        de staart $\sum_{n > N}\abs{u_n}$)}.
  \item (Het recept van Riemann) Zij $t \in \R$. Beschrijf de
        gulzige herschikking van de alternerende harmonische
        \omterm{def:b1:series:def}{reeks}: neem positieve termen $1, \frac13, \frac15, \dots$
        totdat de partiële som $t$ voor het eerst overschrijdt, dan
        negatieve termen $-\frac12, -\frac14, \dots$ totdat zij
        voor het eerst onder $t$ zakt, en herhaal. Toon aan dat
        elke term precies één keer wordt gebruikt, dat de partiële
        sommen na de eerste overschrijding binnen de laatst
        gebruikte term van $t$ blijven, en besluit dat de
        herschikte \omterm{def:b1:series:def}{reeks} naar $t$ convergeert: \emph{elke}
        voorgeschreven som is bereikbaar.
\end{enumerate}

\medskip
\noindent\textbf{Deel V --- De $(p, q)$-formule.} Leg gehele
getallen $p, q \geq 1$ vast. Herschik de alternerende harmonische
\omterm{def:b1:series:def}{reeks} in blokken: $p$ positieve termen (de volgende oneven
omgekeerden), daarna $q$ negatieve termen (de volgende even
omgekeerden), en herhaal.
\begin{enumerate}[resume]
  \item Ga na dat dit een echte herschikking is (elke term precies
        één keer), en dat zij voor $(p, q) = (1, 2)$ luidt
        \[
        1 - \frac12 - \frac14 + \frac13 - \frac16 - \frac18 +
        \frac15 - \frac1{10} - \frac1{12} + \dots
        \]
  \item (De exacte halvering) Bewijs voor $(p, q) = (1, 2)$ de
        blokidentiteit
        \[
        \frac{1}{2k-1} - \frac{1}{4k-2} - \frac{1}{4k}
        = \frac12\Bigl(\frac{1}{2k-1} - \frac{1}{2k}\Bigr),
        \]
        en leid de \emph{exacte} betrekking $T_{3K} = \frac12
        S_{2K}$ af tussen de herschikte partiële sommen en de
        oorspronkelijke: de halvering van de som is in elk eindig
        stadium zichtbaar, niet pas in de limiet.
  \item Toon aan dat de partiële som na $K$ volledige blokken
        gelijk is aan $\sum_{j=1}^{pK} \frac{1}{2j-1} -
        \sum_{j=1}^{qK} \frac{1}{2j}$, en bereken haar limiet met
        vraag 7:
        \[
        \ln 2 + \frac12 \ln\frac pq .
        \]
  \item Beheers de partiële sommen \emph{binnen} een blok (de
        termen naderen tot $0$) en besluit dat de $(p,
        q)$-herschikte \omterm{def:b1:series:def}{reeks} naar $\ln 2 + \frac12\ln \frac pq$
        convergeert. In het bijzonder geeft $(1, 2)$ de waarde
        $\frac{\ln 2}{2}$: toets dit aan de eerste negen termen,
        $T_9 = 0.3083$, die naar $0.3466$ kruipen.
  \item Controles en bereik: $(1,1)$ geeft $\ln 2$ terug; $(2,1)$
        geeft $\frac32\ln 2$; welke sommen zijn met
        $(p, q)$-blokken bereikbaar, en hoe verhoudt dit
        aftelbare menu zich tot de volledige kaart van Riemann
        (vraag 14)?
\end{enumerate}

\medskip
\noindent\textbf{Deel VI --- Epiloog: $\gamma$ aan het werk, en
synthese.}
\begin{enumerate}[resume]
  \item Identificeer de som van de convergente \omterm{def:b1:series:def}{reeks}
        $\sum_{k\geq1} \bigl(\frac1k - \ln\frac{k+1}{k}\bigr)$
        (\cref{exo:b1:series:7}): toon aan dat zij gelijk is aan
        $\gamma$.
  \item Laat het recept van Riemann (vraag 14) lopen voor het doel
        $t = 1$ en som de eerste twaalf voortgebrachte termen op
        ($1, \frac13, -\frac12, \frac15, -\frac14, \frac17,
        \frac19, -\frac16, \frac1{11}, \frac1{13}, -\frac18,
        \frac1{15}$), met berekening van de partiële som
        ($\approx 0.980$) --- kijk hoe het algoritme rond zijn
        doel ademt.
  \item Toon aan dat een zekere herschikking van de alternerende
        harmonische \omterm{def:b1:series:def}{reeks} naar $+\infty$ divergeert \emph{(blokken
        positieve termen die telkens lang genoeg zijn om $1$ te
        winnen, met vraag 12, gescheiden door telkens één
        negatieve term)}.
  \item Scherp \cref{ex:b1:series:harmonicstack} aan met de
        $\gamma$-wet: toon aan dat de eerste index met $H_N \geq
        20$ voldoet aan $N = \eu^{\,20 - \gamma}\,(1 + o(1))
        \approx 2.7\cdot10^{8}$ --- de \omterm{pb:b1:series:1}{constante van Euler} is
        precies de correctie die de grove omsluiting miste.
  \item Synthese, telkens één zin: (i) de $\gamma$-wet en wat elk
        van haar drie stukken ($\ln n$, $\gamma$, $\frac{1}{2n}$)
        bijdraagt; (ii) waarom voorwaardelijke convergentie de som
        van de volgorde laat afhangen terwijl \omterm{thm:b1:series:absolute}{absolute
        convergentie} dat verbiedt; (iii) hoe de
        $(p,q)$-formule een \emph{berekening} met de $\gamma$-wet
        was en geen abstracte bestaansbewering; (iv) waar deze
        draden verdergaan --- machtreeksen en producten van
        \omterm{def:b1:series:def}{reeksen} in het volume van bachelorjaar 2, en de
        weekendopgave van het volume van bachelorjaar 3 over de
        formule van Stirling, waar dezelfde boekhouding van som
        tegenover \omterm{thm:b1:integration:def}{integraal} op volle kracht draait.
\end{enumerate}
\end{problem}
